\chapter{Release 1 : Socle utilisateur et participation citoyenne}
\label{chap:release1}

% Affichage robuste : le rapport reste compilable avant l'export PNG depuis Draw.io.
\providecommand{\releasefigure}[4]{%
\begin{figure}[H]
    \centering
    \IfFileExists{#1}{%
        \includegraphics[width=#2\textwidth]{#1}%
    }{%
        \fbox{\parbox[c][4.2cm][c]{0.88\textwidth}{\centering
        \textbf{Diagramme à exporter depuis Draw.io}\\[0.3cm]
        \texttt{\detokenize{#4}}\\
        vers\\
        \texttt{\detokenize{#1}}}}
    }%
    \caption{#3}
\end{figure}}

\section{Introduction de la Release}
\label{sec:r1_introduction}

Conformément à la planification établie au chapitre~\ref{chap:analyse_specification_besoins}, la première release d'EcoRewind correspond au socle utilisateur et à la participation citoyenne. Elle regroupe deux sprints complémentaires. Le premier sprint est consacré à la gestion des accès, à l'authentification, aux profils et aux rôles. Le deuxième sprint introduit les principales fonctionnalités citoyennes de consultation et d'interaction, en particulier le fil communautaire, les publications, les notifications, la consultation de la carte des points de collecte ainsi que le suivi d'une partie de l'activité de l'utilisateur.

Cette release livre un incrément fonctionnel directement exploitable. Un visiteur peut créer son compte et accéder à la plateforme après vérification. Un utilisateur authentifié peut gérer son profil, consulter des points de collecte, publier du contenu, interagir avec la communauté et recevoir des notifications. Les fonctionnalités plus avancées liées à la collecte intelligente, aux missions terrain, aux équipements IoT, aux contenus éducatifs et aux services de coordination territoriale sont traitées dans les releases suivantes.

\section{Sprint 1 : Authentification, profils et gestion des rôles}
\label{sec:r1_sprint1}

\subsection{Objectifs et backlog du Sprint 1}

L'objectif du Sprint 1 est de mettre en place une gestion sécurisée des identités et des accès. Le système doit permettre la création de compte, la vérification OTP, la connexion par JWT, le renouvellement de session, la récupération du mot de passe, l'authentification sociale, l'activation de la MFA, la consultation du profil et l'administration de base des comptes. Le modèle \texttt{User} centralise les informations d'identité, le rôle, le QR code, le score global, l'avatar ainsi que les attributs liés à la MFA et à la récupération du mot de passe.

\begin{longtable}{|>{\centering\arraybackslash}p{1cm}|p{9.2cm}|>{\centering\arraybackslash}p{1.8cm}|>{\centering\arraybackslash}p{1.8cm}|}
\caption{Backlog du Sprint 1 -- Accès, profils et rôles}\label{tab:r1_s1_backlog}\\
\hline
\rowcolor{gray!35}\textbf{ID} & \textbf{User Story} & \textbf{Priorité} & \textbf{État}\\\hline
\endfirsthead
\hline\rowcolor{gray!35}\textbf{ID} & \textbf{User Story} & \textbf{Priorité} & \textbf{État}\\\hline
\endhead
1.1 & En tant que visiteur, je veux créer et vérifier mon compte afin d'accéder aux services personnalisés de la plateforme. & Élevée & Réalisée\\\hline
1.2 & En tant qu'utilisateur, je veux m'authentifier, renouveler ma session et me déconnecter de manière sécurisée. & Élevée & Réalisée\\\hline
1.3 & En tant qu'utilisateur, je veux consulter et modifier mon profil, mon avatar et mes informations essentielles. & Élevée & Réalisée\\\hline
1.4 & En tant qu'utilisateur, je veux récupérer mon mot de passe et gérer la MFA afin de renforcer la sécurité de mon compte. & Élevée & Réalisée\\\hline
1.5 & En tant qu'utilisateur, je veux me connecter via Google ou Facebook afin de simplifier l'accès à la plateforme. & Moyenne & Réalisée\\\hline
1.6 & En tant qu'administrateur, je veux gérer les comptes et leurs rôles selon mes autorisations. & Élevée & Réalisée\\\hline
\end{longtable}

\subsection{Raffinement des cas d'utilisation}

Le diagramme de cas d'utilisation du Sprint 1 place les acteurs à l'extérieur de la frontière du système EcoRewind. Les cas d'utilisation représentent des objectifs métier orientés utilisateur et non des détails techniques d'implémentation. L'inscription et la vérification du compte sont distinguées comme deux objectifs liés au parcours d'accès. L'authentification constitue une étape préalable aux opérations privées, sans être artificiellement incluse dans tous les cas d'utilisation du profil. La vérification du code MFA étend l'authentification uniquement lorsque cette option est activée. Les opérations de gestion des utilisateurs sont réservées aux rôles administratifs.

\releasefigure{image/usecase_sprint1_access.png}{0.95}{Diagramme de cas d'utilisation du Sprint 1}{usecase_sprint1_access.drawio}
\label{fig:r1_s1_uc}

\subsection{Description textuelle du cas « S'authentifier »}

\begin{longtable}{|p{3.2cm}|p{10.2cm}|}
\caption{Description du cas d'utilisation « S'authentifier »}\label{tab:r1_uc_auth}\\
\hline\rowcolor{gray!20}\textbf{Élément} & \textbf{Description}\\\hline
Acteur principal & Utilisateur disposant d'un compte EcoRewind.\\\hline
Préconditions & Le compte existe, il est actif et l'utilisateur connaît ses identifiants.\\\hline
Postconditions & Un jeton d'accès et un jeton de renouvellement sont produits, sauf si une vérification MFA reste requise. Les informations minimales de session sont également retournées à l'application Flutter.\\\hline
Scénario nominal & 1. L'utilisateur saisit son courriel et son mot de passe. \newline 2. Flutter transmet les données à \texttt{POST /token}. \newline 3. Le backend recherche l'utilisateur puis vérifie le mot de passe haché. \newline 4. Si la MFA n'est pas activée, le backend génère un jeton d'accès et un jeton de renouvellement. \newline 5. Il synchronise les données nécessaires avec Firebase. \newline 6. Il retourne les jetons, le rôle, l'identifiant, l'avatar et le QR personnel. \newline 7. Flutter ouvre l'espace correspondant au rôle.\\\hline
Alternatives & A1 : identifiants invalides, réponse HTTP 401. \newline A2 : MFA active, retour d'un \texttt{mfa\_token}, saisie du code TOTP puis appel de \texttt{/auth/mfa/verify}. \newline A3 : compte désactivé, accès refusé.\\\hline
\end{longtable}

\subsection{Diagrammes de séquence}

Le diagramme de la figure~\ref{fig:r1_s1_seq_auth} respecte le flux réel de \texttt{POST /token}. Les objets impliqués représentent le niveau métier et applicatif pertinent : acteur, application Flutter, API FastAPI, base de données et service Firebase. Le fragment \texttt{alt} distingue l'échec d'authentification, la connexion directe et la branche MFA. Les retours sont représentés en pointillés et les appels synchrones en traits pleins.

\releasefigure{image/sequence_sprint1_authentication.png}{0.98}{Diagramme de séquence « S'authentifier »}{sequence_sprint1_authentication.drawio}
\label{fig:r1_s1_seq_auth}

La création du compte suit une logique en deux temps : inscription puis vérification. Le diagramme de la figure~\ref{fig:r1_s1_seq_signup} montre la vérification de l'unicité du courriel, le hachage du mot de passe, la persistance de l'utilisateur, l'initialisation des données utiles au QR code et à Firebase, puis l'envoi du code OTP.

\releasefigure{image/sequence_sprint1_signup.png}{0.98}{Diagramme de séquence « Créer un compte et vérifier son accès »}{sequence_sprint1_signup.drawio}
\label{fig:r1_s1_seq_signup}

\subsection{Diagramme de classes du Sprint 1 — CORRIGÉ}
\label{subsec:r1_s1_class_corrected}

Le diagramme de classes du Sprint 1 a été validé et corrigé pour refléter fidèlement le code source réel. Les corrections majeures apportées sont :

\begin{itemize}
    \item \textbf{Classe User :} 15 attributs détaillés, y compris \texttt{global\_score} (suivi de l'impact), \texttt{fcm\_token} (notifications Firebase), \texttt{mfa\_enabled} et \texttt{mfa\_secret} (authentification forte TOTP).
    \item \textbf{Classe OTPCode :} 6 attributs complets incluant \texttt{purpose} (register, reset), \texttt{expires\_at} (validité du code) et \texttt{is\_used} (statut d'utilisation).
    \item \textbf{Annotations UML :} Notation stricte des clés primaires (PK), étrangères (FK) et uniques (UK). Les champs nullable sont explicitement marqués.
    \item \textbf{Méthodes :} La méthode \texttt{generate\_unique\_qr\_token()} est documentée ; elle génère un identifiant unique au format \texttt{ECOREWIND-UUID4}.
    \item \textbf{Absence de relations explicites :} Le diagramme de Sprint 1 ne montre pas les relations vers Post, Comment, Notification, etc., car ces entités appartiennent au Sprint 2.
\end{itemize}

Le diagramme est accessible en format Draw.io modifiable à l'adresse :
\begin{center}
    \texttt{uml\_diagrams/sprint1\_class\_diagram.drawio}
\end{center}

\releasefigure{image/class_sprint1_access.png}{0.90}{Diagramme de classes du Sprint 1 (CORRIGÉ)}{class_sprint1_access.drawio}
\label{fig:r1_s1_class}

\subsection{Scénarios de validation}

Les scénarios de validation du Sprint 1 portent sur l'inscription avec courriel unique, le refus d'identifiants invalides, la production des JWT, la branche MFA, la récupération de mot de passe, la protection d'une ressource privée et le contrôle des rôles d'administration. Ils constituent des critères de recette métier ; ils ne remplacent pas l'exécution des tests automatisés mais en reflètent les objectifs fonctionnels.

\section{Sprint 2 : Participation citoyenne, publications, notifications et consultation des points}
\label{sec:r1_sprint2}

\subsection{Objectifs et backlog du Sprint 2}

Le Sprint 2 introduit les principales fonctionnalités citoyennes prévues dans la première release. Il permet à l'utilisateur de consulter le fil d'actualité, de publier du contenu, d'interagir avec les publications, de recevoir des notifications, de consulter les points de collecte sur carte et d'accéder à une première lecture de son activité dans la plateforme. Le dépôt montre également que les publications peuvent être modérées automatiquement par un pipeline IA prudent avant validation finale.

\begin{longtable}{|>{\centering\arraybackslash}p{1cm}|p{9.2cm}|>{\centering\arraybackslash}p{1.8cm}|>{\centering\arraybackslash}p{1.8cm}|}
\caption{Backlog du Sprint 2 -- Participation citoyenne et consultation}\label{tab:r1_s2_backlog}\\
\hline\rowcolor{gray!35}\textbf{ID} & \textbf{User Story} & \textbf{Priorité} & \textbf{État}\\\hline
\endfirsthead
\hline\rowcolor{gray!35}\textbf{ID} & \textbf{User Story} & \textbf{Priorité} & \textbf{État}\\\hline
\endhead
2.1 & En tant que citoyen, je veux consulter le fil d'actualité afin de suivre les actions partagées par la communauté. & Élevée & Réalisée\\\hline
2.2 & En tant que citoyen, je veux publier une action citoyenne avec un média éventuel afin de partager mon engagement écologique. & Élevée & Réalisée\\\hline
2.3 & En tant que citoyen, je veux aimer, commenter, répondre et enregistrer une publication afin d'interagir avec la communauté. & Élevée & Réalisée\\\hline
2.4 & En tant qu'utilisateur, je veux recevoir des notifications afin d'être informé des interactions liées à mon activité. & Élevée & Réalisée\\\hline
2.5 & En tant que visiteur ou citoyen, je veux consulter les points de collecte sur une carte afin de localiser les centres disponibles. & Élevée & Réalisée\\\hline
2.6 & En tant que citoyen, je veux rechercher et filtrer les points puis consulter leurs détails. & Élevée & Réalisée\\\hline
2.7 & En tant que citoyen, je veux soumettre un témoignage ou proposer un point de collecte afin de participer à l'amélioration du service. & Moyenne & Réalisée\\\hline
2.8 & En tant qu'utilisateur, je veux consulter certaines statistiques personnelles d'activité et d'impact. & Moyenne & Réalisée\\\hline
\end{longtable}

\subsection{Raffinement des cas d'utilisation}

Le diagramme de cas d'utilisation du Sprint 2 met en évidence la participation citoyenne autour de trois axes : la consultation, l'interaction et la contribution. La consultation du fil et de la carte constitue la base des usages publics ou authentifiés. L'ajout d'un média étend la création d'une publication car il reste facultatif. Les interactions telles que le commentaire, la réponse, le J'aime et la sauvegarde sont déclenchées à la demande et ne doivent pas être confondues avec des détails techniques. La modération automatique fait partie du traitement interne de la publication, tandis que les notifications sont déclenchées en réaction aux interactions sociales.

\releasefigure{image/usecase_sprint2_participation.png}{0.97}{Diagramme de cas d'utilisation du Sprint 2}{usecase_sprint2_participation.drawio}
\label{fig:r1_s2_uc}

\subsection{Description textuelle du cas « Publier une action citoyenne »}

\begin{longtable}{|p{3.2cm}|p{10.2cm}|}
\caption{Description du cas d'utilisation « Publier une action citoyenne »}\label{tab:r1_uc_post}\\
\hline\rowcolor{gray!20}\textbf{Élément} & \textbf{Description}\\\hline
Acteur principal & Citoyen authentifié.\\\hline
Préconditions & Le jeton d'accès est valide ; les données envoyées respectent le schéma attendu ; si un média est joint, il est d'un format accepté.\\\hline
Postconditions & La publication est persistée avec le statut \texttt{pending\_review} et une tâche de modération est planifiée. Le statut final peut évoluer vers \texttt{published} ou rester en attente de validation humaine.\\\hline
Scénario nominal & 1. Le citoyen rédige son contenu et ajoute éventuellement un média. \newline 2. Flutter téléverse le média si nécessaire et récupère son URL. \newline 3. Flutter appelle \texttt{POST /posts}. \newline 4. L'API persiste immédiatement la publication avec le statut \texttt{pending\_review}. \newline 5. Elle répond à l'application sans attendre la fin de l'analyse. \newline 6. Une tâche de fond exécute la modération IA. \newline 7. Le statut est mis à jour et une notification informe l'auteur.\\\hline
Alternatives & A1 : média invalide ou trop volumineux, la requête d'upload est rejetée. \newline A2 : contenu jugé sûr, statut \texttt{published}. \newline A3 : contenu signalé ou erreur du pipeline, maintien en \texttt{pending\_review} pour arbitrage humain.\\\hline
\end{longtable}

\subsection{Diagrammes de séquence}

Le Sprint 2 mobilise deux diagrammes de séquence complémentaires. Le premier décrit la consultation des points de collecte depuis la carte. Il ne contient pas de fragment \texttt{ref S'authentifier}, car cette fonctionnalité reste publique dans sa forme de base. Le fragment \texttt{alt} distingue les cas classiques : liste disponible, liste vide ou erreur technique.

\releasefigure{image/sequence_sprint2_collection_points.png}{0.98}{Diagramme de séquence « Consulter les points de collecte »}{sequence_sprint2_collection_points.drawio}
\label{fig:r1_s2_seq_points}

Le second diagramme décrit la création d'une publication citoyenne. Le fragment \texttt{opt [média fourni]} précède la création. La réponse HTTP est renvoyée avant la fin de la modération, ce qui reflète correctement le caractère asynchrone du traitement. Le fragment \texttt{alt} distingue l'approbation automatique et la révision humaine.

\releasefigure{image/sequence_sprint2_publish_post.png}{0.99}{Diagramme de séquence « Publier une action citoyenne »}{sequence_sprint2_publish_post.drawio}
\label{fig:r1_s2_seq_post}

\subsection{Diagramme de classes du Sprint 2 — COMPLET ET CORRIGÉ}
\label{subsec:r1_s2_class_corrected}

Le diagramme de classes du Sprint 2 a été entièrement validé et corrigé pour inclure toutes les entités réelles du code source. Les corrections et améliorations apportées sont :

\begin{itemize}
    \item \textbf{Classe Post :} 14 attributs complets incluant les champs de modération IA (\texttt{status}, \texttt{moderation\_score}, \texttt{moderation\_reason}, \texttt{moderation\_details}, \texttt{moderated\_at}, \texttt{moderation\_model\_version}).
    
    \item \textbf{Classes d'association :} \texttt{Like} et \texttt{SavedPost} sont correctement modélisées comme des classes d'association (many-to-many) avec leurs propres attributs et clés primaires.
    
    \item \textbf{Classe Comment :} 8 attributs avec la relation réflexive \texttt{parent\_id} permettant la modélisation des réponses aux commentaires (\textit{comment replies}).
    
    \item \textbf{Classe Notification :} 12 attributs détaillant les types de notifications (like, comment, save, intercommunality\_message, actor\_reply), les références croisées (\texttt{post\_id}, \texttt{comment\_id}, \texttt{source\_notification\_id}) et le statut de lecture (\texttt{is\_read}).
    
    \item \textbf{Classe Testimonial :} 8 attributs permettant aux utilisateurs de soumettre des retours et témoignages avec système d'approbation admin (\texttt{is\_approved}) et de mise en avant (\texttt{is\_featured}).
    
    \item \textbf{Classe CenterProposal :} 10 attributs permettant aux citoyens de proposer de nouveaux points de collecte avec suivi du statut (\texttt{pending}, \texttt{approved}, \texttt{rejected}).
    
    \item \textbf{Classe CollectionPoint :} 10 attributs représentant les points de collecte consultables avec suivi de l'état de charge (\texttt{load\_level}), du statut opérationnel et des types de déchets acceptés.
    
    \item \textbf{Multiplicité des associations :} Notation stricte UML 2.5 (1..*, 0..1, etc.) reflétant les relations réelles dans la base de données.
\end{itemize}

Le diagramme Sprint 2 étend logiquement le Sprint 1 en réutilisant la classe \texttt{User} et en ajoutant tous les domaines fonctionnels de la participation citoyenne. Les fichiers sont accessibles en format Draw.io modifiable à :
\begin{center}
    \texttt{uml\_diagrams/sprint2\_class\_diagram.drawio}
\end{center}

\releasefigure{image/class_sprint2_participation.drawio.png}{0.98}{Diagramme de classes du Sprint 2 (COMPLET ET CORRIGÉ)}{class_sprint2_participation.drawio}
\label{fig:r1_s2_class}

\subsection{Scénarios de validation}

Les scénarios de validation du Sprint 2 portent sur :
\begin{itemize}
    \item Consultation du fil d'actualité avec pagination
    \item Création d'une publication avec/sans média
    \item Upload et validation des formats de médias
    \item Modération automatique prudente (pipeline IA)
    \item Ajout de commentaires et réponses aux commentaires
    \item Interactions (J'aime, sauvegarde)
    \item Réception des notifications suite aux interactions
    \item Consultation de la carte avec filtrage/recherche
    \item Affichage des points de collecte avec détails
    \item Soumission de témoignages et propositions de centres
\end{itemize}

Ces scénarios vérifient que la participation citoyenne et l'accès à l'information fonctionnent de manière cohérente dans l'incrément livré.

\section{Implémentation de la Release 1}
\label{sec:r1_implementation}

L'application cliente est développée avec Flutter pour les environnements mobiles et web. Elle communique avec les routeurs FastAPI par HTTP, gère les jetons de session côté client et adapte l'interface au rôle courant. Le backend repose sur une organisation modulaire autour des domaines \texttt{app/auth}, \texttt{app/users}, \texttt{app/posts}, \texttt{app/notifications}, \texttt{app/collection\_points} et \texttt{app/community}. Les données métier sont manipulées via SQLAlchemy et l'évolution du schéma est assurée par Alembic.

La première release mobilise plus particulièrement les modules d'accès, de participation citoyenne, de consultation cartographique et de notifications. Firebase complète le backend pour certains usages transverses, notamment la synchronisation de données liées au profil, au score ou aux notifications. Les mécanismes de modération IA présents dans le dépôt enrichissent le fil communautaire tout en conservant une décision humaine sur les cas sensibles.

\subsection{Architecture modulaire du backend}

L'architecture feature-based du backend est organisée comme suit :

\begin{itemize}
    \item \textbf{app/auth/} — Gestion de l'authentification, OTP, JWT, MFA (TOTP), réinitialisation du mot de passe
    \item \textbf{app/users/} — Modèle User, gestion des profils, génération du QR code unique
    \item \textbf{app/posts/} — Modèles Post, Comment, Like, SavedPost ; gestion du fil communautaire
    \item \textbf{app/notifications/} — Modèle Notification ; déclenchement asynchrone des alertes utilisateur
    \item \textbf{app/collection\_points/} — Modèle CollectionPoint ; consultation cartographique et filtrage
    \item \textbf{app/community/} — Modèles Testimonial, CenterProposal ; contributions citoyennes
\end{itemize}

Chaque module contient :
\begin{itemize}
    \item \textbf{models.py} — Classes SQLAlchemy définissant la structure des tables
    \item \textbf{schemas.py} — Schémas Pydantic pour validation/sérialisation des requêtes/réponses
    \item \textbf{router.py} — Endpoints FastAPI exposés au client
    \item \textbf{repository.py} (optionnel) — Couche d'accès aux données
    \item \textbf{service.py} (optionnel) — Logique métier transversale
\end{itemize}

\subsection{Technologies clés de Release 1}

\begin{itemize}
    \item \textbf{FastAPI 0.95+} — Framework web performant avec validation automatique, documentation OpenAPI
    \item \textbf{SQLAlchemy 2.0+} — ORM avec support complet des relations complexes (many-to-many, reflexive)
    \item \textbf{Pydantic v2} — Validation de schémas avec mécanismes de sérialisation flexible
    \item \textbf{Alembic} — Gestion des migrations de schéma de base de données
    \item \textbf{PyJWT} — Génération et validation des jetons JWT pour l'authentification stateless
    \item \textbf{python-jose} — Cryptage TOTP et gestion MFA
    \item \textbf{Firebase Admin SDK} — Synchronisation des données critiques et notifications push
    \item \textbf{Flutter} — Client mobile/web unifiés avec gestion de session côté client
\end{itemize}

\section{Accès aux diagrammes corrigés}
\label{sec:r1_diagrams_access}

Tous les diagrammes UML ont été corrigés et validés contre le code source réel d'EcoRewind. Les fichiers sont disponibles en format Draw.io (modifiable) dans le répertoire du projet :

\begin{center}
    \texttt{uml\_diagrams/}
\end{center}

\noindent\textbf{Fichiers disponibles :}
\begin{itemize}
    \item \texttt{sprint1\_class\_diagram.drawio} — Diagramme de classes Sprint 1 (User, OTPCode)
    \item \texttt{sprint2\_class\_diagram.drawio} — Diagramme de classes Sprint 2 (Post, Comment, Like, SavedPost, Notification, Testimonial, CenterProposal, CollectionPoint)
    \item \texttt{DOWNLOAD\_DIAGRAMS.html} — Page HTML interactive pour télécharger ou ouvrir les diagrammes
\end{itemize}

\noindent\textbf{Comment accéder aux diagrammes :}
\begin{enumerate}
    \item Ouvrir le fichier HTML : \texttt{uml\_diagrams/DOWNLOAD\_DIAGRAMS.html}
    \item Cliquer sur le bouton « Télécharger » pour obtenir le fichier .drawio
    \item Importer le fichier dans \url{https://app.diagrams.net} pour édition ou visualisation
\end{enumerate}

\section{Conclusion de la Release 1}
\label{sec:r1_conclusion}

Cette release fournit un incrément cohérent avec la planification établie dans le chapitre d'analyse et avec les fonctionnalités réellement observées dans le dépôt de code. Elle installe le socle d'accès sécurisé avec authentification multi-facteurs, permet la participation citoyenne à travers le fil communautaire enrichi de mécanismes de modération IA, et offre une première consultation des points de collecte avec contributions citoyennes.

Les diagrammes UML associés, désormais corrigés et validés, décrivent fidèlement :
\begin{itemize}
    \item Les responsabilités fonctionnelles réellement implémentées
    \item Les échanges entre Flutter, FastAPI et les services de support (Firebase, PostgreSQL)
    \item Les principales entités métier et leurs relations conformes à la notation UML 2.5
\end{itemize}

La Release 1 prépare ainsi l'intégration des volets plus avancés du projet, notamment la collecte intelligente avec équipements IoT, les missions terrain pour les collecteurs, le suivi des équipements, la sensibilisation éducative et les mécanismes de coordination territoriale prévus dans les releases suivantes.
